\chapter{What is Risk in a Construction and Trades Context?}
\label{chap:What is Risk in a Construction and Trades Context?}

%---------------------------- SECTION ----------------------------
\section{Why this assessment matters in construction}

Construction and trades work sits among the highest-risk categories of employment in the United Kingdom. The Health and Safety Executive consistently records more fatal injuries and major injuries per 100,000 workers in construction than in almost any other sector. The built environment is created by people who work at height, in excavations, alongside heavy plant, within confined spaces, and in the presence of hazardous substances --- often simultaneously, on the same site, within metres of one another. Understanding what risk actually means in this context is therefore not an academic exercise; it is the foundation upon which every decision a site manager, principal contractor, or sole trader makes must rest.

Risk, in the technical sense used by health and safety law and practice, is a function of two components: the likelihood that a harm will occur, and the severity of that harm if it does. A fragile roof is not inherently a risk unless someone is likely to access it; a box of caustic cleaning fluid is not inherently a risk unless it can be opened, spilled, or misused. Construction amplifies both dimensions. The dynamic nature of sites means that likelihood of exposure to hazards changes daily as structures rise, ground conditions alter, and different trades interact. The severity dimension is equally pronounced: falls from height kill, amputations are irreversible, occupational lung disease caused by silica or asbestos is progressive and incurable. The risk equation in construction is therefore not theoretical --- the consequences of miscalculation are measured in lives, livelihoods, and permanent impairment.

A further complexity unique to construction is the transient and multi-organisational nature of the workforce. Unlike a factory where the same operatives perform the same tasks at the same workstation each day, a construction site brings together dozens of contractors, subcontractors, and self-employed trades, each arriving at different phases, each carrying their own risk profiles, each potentially interacting with hazards created or modified by others. A groundworker's excavation becomes the structural engineer's concern; a painter's solvent fumes become the electrician's ignition risk. Risk, in this environment, is cumulative, interactive, and often invisible to those who have not been trained to look for it.

For these reasons, the very first intellectual step that any competent person must take before committing resources, writing a method statement, or issuing an induction pack is to understand the risk landscape of the specific work being undertaken. This chapter establishes that conceptual foundation, defining risk in practical terms, describing the principal hazard categories present across construction and trades, and explaining the legislative duty to assess, record, and manage risk as a prerequisite for every subsequent chapter in this volume.

\barquote{``Those who cannot remember the past are condemned to repeat it.'' The same is true of hazards that are known but unassessed --- they will eventually result in harm.}

\example{Site Reality}{A residential refurbishment project in an urban terraced row employs a main contractor and six subcontractors concurrently: a groundworker laying drainage, an electrician first-fixing throughout, a plasterer working on the first floor, a roofer replacing tiles, a scaffold firm maintaining access, and a window installer. Each trade has its own risk profile. The groundworker's open trench creates a fall risk for every other operative on site. The roofer's tile removal creates falling-object risk for the scaffold operative below. The electrician's trailing cables on unfinished floors create a trip hazard for the plasterer carrying loaded boards. Without a site-wide risk assessment that maps these interactions, each trade assesses only its own work, and the interfaces --- where most incidents occur --- remain unmanaged.}

%---------------------------- SECTION ----------------------------
\section{Legal and professional framework}

The legal duty to assess and manage risk in UK construction derives from a hierarchy of primary and secondary legislation, supported by Approved Codes of Practice and authoritative HSE guidance. The Health and Safety at Work etc.\ Act 1974 establishes the overarching duty on employers and the self-employed to ensure, so far as is reasonably practicable, the health, safety, and welfare of employees and those affected by their undertaking. This ``so far as is reasonably practicable'' standard requires a proportionate weighing of the cost and effort of risk reduction against the magnitude of the risk; in construction, where consequences are frequently catastrophic, this threshold is regularly interpreted as demanding substantial and demonstrable control measures.

The Management of Health and Safety at Work Regulations 1999 operationalise the general duty by imposing an explicit requirement to carry out suitable and sufficient risk assessments, implement preventive and protective measures, appoint competent persons, and establish procedures for serious and imminent danger. The Construction (Design and Management) Regulations 2015 add a project-specific layer, distributing duties across clients, principal designers, principal contractors, contractors, and workers, and requiring that health and safety be embedded in the design and planning process from inception. HSG65, the HSE's managing for health and safety framework, provides the Plan--Do--Check--Act model that underpins proportionate, continuous risk management across all scales of construction activity.

\paragraph{Key frameworks relevant to this chapter include: the Health and Safety at Work etc.\ Act 1974; the Management of Health and Safety at Work Regulations 1999; the Construction (Design and Management) Regulations 2015; and HSG65.}

\begin{itemize}
  \bitem{Health and Safety at Work etc.\ Act 1974}{The primary enabling legislation placing general duties on employers, the self-employed, and those in control of premises to protect employees and non-employees from risks arising from work activities. Section~2 and Section~3 are the cornerstones of all construction health and safety obligations.}
  \bitem{Management of Health and Safety at Work Regulations 1999}{Requires every employer to make a suitable and sufficient assessment of risks to employees and non-employees, to implement preventive measures on the basis of the general principles of prevention, and to ensure that assessments are reviewed when no longer valid.}
  \bitem{Construction (Design and Management) Regulations 2015}{Applies to all construction projects, distributing duties to clients, principal designers, principal contractors, contractors, and workers. Requires pre-construction information, the health and safety file, and the construction phase plan to be produced and maintained.}
  \bitem{HSG65 --- Managing for Health and Safety}{HSE's primary guidance framework providing the Plan--Do--Check--Act model for proportionate health and safety management, including risk profiling, implementation of controls, monitoring, and review.}
\end{itemize}

%---------------------------- SECTION ----------------------------
\section{Conducting the assessment using the five-step method}

\subsection{Step 1 --- Identify hazards}
\paragraph{In a construction and trades context, hazard identification must look beyond the obvious physical dangers to encompass the full spectrum of hazards present at each phase of a project. Physical hazards include falls from height, being struck by moving vehicles or objects, contact with mobile plant, structural collapse, and manual handling. Health hazards include inhalation of respirable crystalline silica (RCS) from cutting masonry, concrete, or stone; asbestos fibres from disturbing legacy materials in pre-2000 buildings; noise above the lower and upper exposure action values from power tools, breakers, and plant; and hand-arm vibration from prolonged use of vibrating tools. Environmental hazards include contaminated land, buried services, and unstable ground. Hazard identification must be site-specific, phase-specific, and trade-specific; a generic library list of hazards is not a substitute for a walk-through assessment conducted by a competent person who knows the site, the work sequence, and the workforce.}

\subsection{Step 2 --- Decide who or what is affected}
\paragraph{Construction sites contain multiple categories of person who may be harmed, and the assessment must address each. Employees of the principal contractor and each subcontractor are the primary population, but self-employed operatives, labour-only subcontractors, and apprentices or trainees --- who may be less experienced and more vulnerable --- must be specifically considered. Visitors, including clients, designers, inspectors, and delivery drivers, are frequently present and may not be briefed on site hazards. Members of the public adjacent to site boundaries, including pedestrians, cyclists, and residents, represent a significant exposure group, particularly in urban environments where hoardings, pavement licences, and crane oversails bring construction activity into close proximity with uninvolved people. The assessment must also consider workers with particular vulnerabilities, including new and expectant mothers, young workers, and those with health conditions that may increase their susceptibility to specific hazards.}

\subsection{Step 3 --- Evaluate likelihood and consequence}
\paragraph{Risk evaluation in construction should use a structured matrix approach, recording both the inherent risk (the level of risk before any controls are applied) and the residual risk (the level remaining after controls are in place). Likelihood scoring should reflect realistic site conditions --- how often will workers be exposed to the hazard, for how long, and how many will be exposed simultaneously? Consequence scoring must reflect worst-case outcomes; a fall from scaffold does not result in a sprained ankle; it results in death or serious permanent injury. Common risk matrices use a five-by-five or three-by-three scale, producing a numerical risk rating that determines the priority and nature of the controls required. Critically, the evaluation must be conducted by a competent person with sufficient knowledge of the work to make realistic judgements, not by someone who has never been on site or who does not understand the task.}

\subsection{Step 4 --- Select and implement controls}
\paragraph{Controls in construction must be selected by applying the hierarchy of controls in sequence: elimination first, then substitution, then engineering controls, then administrative controls, and finally personal protective equipment. A control is not adequate merely because it has been written down; it must be physically implemented, communicated to those affected, and verified as effective. Each control measure must have a named owner responsible for its implementation and maintenance, and a realistic timescale must be established for measures that cannot be implemented immediately. Method statements, permits to work, toolbox talks, safety inductions, site rules, and PPE provision are all controls, but they must be proportionate to the risk level established at Step 3. PPE must never be the sole or primary control for a high-residual-risk activity.}

\subsection{Step 5 --- Record and review}
\paragraph{Employers with five or more employees are legally required to record the significant findings of their risk assessments in writing. Good practice --- and the expectation of principal contractors, insurers, and enforcement authorities --- is that all significant risk assessments are documented regardless of workforce size. The record must include the hazards identified, the persons at risk, the controls in place, the residual risk rating, the review date, and the name of the competent person who conducted the assessment. Reviews must be triggered not only by the passage of time but by any change to the work activity, method, location, or workforce, or by the occurrence of any near-miss, incident, or enforcement action. On construction sites, where conditions change daily, this means that formal review processes must be integrated into the site management routine.}

\example{Five-Step Application}{A brick-cutting operation on a housing development is assessed using the five-step method. Step 1 identifies RCS as the primary health hazard, noise as a secondary hazard, and the risk of blade contact as a physical hazard. Step 2 identifies the bricklayer, adjacent trades within ten metres, and the site manager conducting inspections as those at risk. Step 3 evaluates the inherent risk of RCS as high (frequent, prolonged exposure; consequence is incurable silicosis and lung cancer). Step 4 implements on-tool water suppression as the primary engineering control, restricts cutting to designated areas away from other trades, limits daily exposure through task rotation, and provides FFP3 respiratory protective equipment as a supplementary last-resort control. Step 5 records the assessment, assigns the site manager as review owner, and sets a review trigger at any change to the cutting method or equipment.}

%---------------------------- SECTION ----------------------------
\section{Applying the hierarchy of controls}

\paragraph{The hierarchy of controls is the organising principle of risk management in UK construction and is embedded in the general principles of prevention set out in Schedule 1 to the Management of Health and Safety at Work Regulations 1999. It requires that every effort to reduce risk to the lowest reasonably practicable level be made by progressing through the hierarchy in sequence, applying higher-order controls wherever feasible before resorting to lower-order measures. In construction, where the temptation to issue PPE and consider the matter resolved is ever-present, the hierarchy demands a more rigorous and creative approach.}

\begin{itemize}
  \bitem{Elimination}{The most effective control is to remove the hazard entirely. In construction, this means designing out risks at the planning stage: specifying pre-cast concrete elements rather than in-situ reinforced concrete to eliminate formwork and working at height; selecting low-silica aggregate in concrete mixes; choosing pre-decorated components to eliminate on-site painting at height; relocating plant routes to remove pedestrian-plant interfaces. CDM 2015 places explicit duties on designers to eliminate foreseeable risks during the design phase.}
  \bitem{Substitution}{Where elimination is not reasonably practicable, the hazardous process, substance, or equipment should be replaced with a less hazardous alternative. This includes substituting wet methods for dry cutting to reduce RCS; using battery-powered tools with lower vibration magnitudes in place of compressed-air equivalents; replacing solvent-based adhesives with water-based alternatives; and specifying lower-silica brick and block types where structural requirements permit.}
  \bitem{Engineering controls}{Physical barriers, guarding, extraction systems, and mechanical aids that reduce exposure without relying on individual behaviour. Examples include fixed edge protection and collective fall arrest systems; on-tool local exhaust ventilation (LEV) or water suppression on cutting and grinding equipment; mechanical lifting aids to eliminate manual handling of heavy materials; acoustic enclosures or barriers around high-noise operations; and interlocked machine guarding on static plant.}
  \bitem{Administrative controls}{Procedures, systems of work, training, supervision, and communication measures that control risk through human action and organisation. In construction these include permit-to-work systems for confined spaces, hot work, and live electrical work; toolbox talks and pre-task briefings; site induction and emergency procedures; health surveillance for workers exposed to noise, vibration, and COSHH substances; task rotation to limit cumulative exposure to vibration; and segregated pedestrian and plant routes.}
  \bitem{PPE}{The last resort in the hierarchy, to be used only where residual risk remains after all higher-order controls have been applied. Construction PPE includes hard hats, safety footwear, high-visibility clothing, gloves, eye protection, hearing protection, and respiratory protective equipment. PPE must be selected for the specific hazard, correctly fitted, maintained in good condition, and worn consistently; a FFP3 respirator that is stored in a van is not a control.}
\end{itemize}

%---------------------------- SECTION ----------------------------
\section{Cadence of assessment and review triggers}

\paragraph{Risk assessments for construction activities are not static documents produced once at tender stage and filed until a HSE visit. The dynamic nature of sites means that assessments must be reviewed at defined intervals and re-examined whenever the conditions that gave rise to the original assessment change. The minimum review cadence should be established during the construction phase planning process, with more frequent reviews assigned to higher-risk activities and more stable review cycles applied to lower-risk ancillary operations.}

\begin{itemize}
  \bitem{Project phase change}{Any transition between major project phases --- from groundworks to superstructure, from structure to fit-out, from fit-out to commissioning --- requires a comprehensive review of all current risk assessments, as the hazard profile of the site changes fundamentally at each phase.}
  \bitem{Change to method or process}{If the planned method of work is altered for any reason --- change in design, substitution of plant or equipment, change in material specification --- the risk assessment for the affected activity must be reviewed and updated before work recommences under the new method.}
  \bitem{Change in workforce}{The arrival of a new subcontractor, the replacement of key competent persons, the engagement of young workers or trainees, or significant changes in gang composition require review of assessments and targeted induction to ensure that new operatives are briefed on current controls.}
  \bitem{Incident, near-miss, or enforcement action}{Any recordable incident under RIDDOR 2013, any near-miss reported through the site safety management system, or any enforcement notice or improvement notice from HSE or a local authority inspector must trigger an immediate review of the relevant risk assessment.}
  \bitem{Changes in legislation, guidance, or standards}{The publication of new or revised Approved Codes of Practice, HSE guidance, British Standards, or other authoritative guidance relevant to the activities being assessed must trigger a review to ensure that controls remain consistent with current best practice and legal compliance.}
\end{itemize}

%---------------------------- SECTION ----------------------------
\section{Common failure modes}

\begin{itemize}
  \bitem{Generic or template assessments}{The single most common failure mode in construction risk assessment is the use of unmodified generic documents. A risk assessment downloaded from a trade association website and applied without modification to a specific site, specific workforce, and specific method of work is not suitable and sufficient. It will not reflect actual site conditions, actual hazard interactions, or actual controls implemented.}
  \bitem{Absence of inherent risk recording}{Many assessments record only the residual risk after controls, without documenting the inherent risk before controls. This makes it impossible to verify whether the controls applied are proportionate to the risk, and prevents meaningful comparison over time as controls are added or removed.}
  \bitem{Controls without named owners}{Listing a control measure without assigning a specific named individual responsible for its implementation, maintenance, and verification is a structural failure. Controls that belong to everyone belong to no one; on a site with multiple contractors, this produces gaps that are exploited by circumstance.}
  \bitem{Assessments inaccessible to the workforce}{Risk assessments written in dense technical language, kept in an office filing system, or produced only in English on sites with multilingual workforces are functionally useless as tools for protecting workers. The assessment must be communicated to those it protects in a form they can understand and act upon.}
  \bitem{Failure to assess health hazards}{Physical hazards dominate most construction risk assessments, but occupational health hazards --- RCS, asbestos, noise, vibration, COSHH substances, and manual handling --- cause the majority of work-related ill-health in the sector. Assessments that omit or superficially treat these hazards are structurally incomplete.}
  \bitem{Review cycle not implemented}{Risk assessments with review dates that have passed without action are a recurring finding in HSE investigations following fatal accidents. The review schedule must be actively managed; it is not sufficient to print a review date on a document if no organisational process exists to ensure that review actually happens.}
\end{itemize}

\example{Failure Pattern}{A groundworks subcontractor submits a risk assessment for trench excavation on an inner-city mixed-use development. The document is a standard template used across all their projects; the site address has been updated but no other modification has been made. The assessment does not reference the gas main confirmed at 900\,mm depth by the utility survey, does not address the proximity of the public footpath on the northern boundary, and does not record any controls for the water main at 1,200\,mm depth that requires the trench to pass within 300\,mm. The assessment shows a residual risk rating of ``low'' but records no inherent risk. When a breaker strikes the gas main during open-cut excavation, the resulting incident triggers a HSE investigation that identifies the inadequate risk assessment as a primary causal factor.}

%---------------------------- CHAPTER SUMMARY ----------------------------
\section{Chapter Summary}

\begin{table}[H]
\centering
\begin{tabularx}{\textwidth}{>{\raggedright\arraybackslash}p{3.2cm}X}
\toprule
\textbf{Assessment Element} & \textbf{Operational Requirement} \\
\midrule
Legal/Professional Basis & HSWA 1974 s.2--3; Management Regulations 1999 reg.3; CDM 2015; HSG65. Risk assessment is a legal duty, not a discretionary management tool. \\
Method & Apply the full five-step process and document inherent/residual risk. \\
Controls & Apply the hierarchy of controls in sequence; elimination and engineering controls must precede PPE; each control requires a named owner. \\
Cadence & Review at each phase change, method change, workforce change, incident, or regulatory update; minimum six-monthly review cycle for ongoing activities. \\
Assurance & Use audit, incident learning, and governance reporting to verify effectiveness. \\
\bottomrule
\end{tabularx}
\end{table}

\begin{itemize}
  \bitem{Key takeaway 1}{Risk in construction is a function of likelihood and consequence; both dimensions are elevated in this sector, making structured assessment and management a non-negotiable operational requirement rather than a compliance exercise.}
  \bitem{Key takeaway 2}{The five-step risk assessment method must be applied site-specifically, trade-specifically, and phase-specifically; generic template documents do not constitute suitable and sufficient assessments under the Management Regulations 1999.}
  \bitem{Key takeaway 3}{The hierarchy of controls demands that elimination, substitution, and engineering controls are exhausted before administrative controls and PPE are relied upon; this sequence is embedded in Schedule 1 of the Management Regulations 1999 as a legal requirement.}
  \bitem{Key takeaway 4}{Occupational health hazards --- silica, asbestos, noise, vibration, and manual handling --- cause the majority of work-related ill-health in construction; risk assessments that address only physical safety hazards are structurally and legally incomplete.}
\end{itemize}

\barquote{In Chapter~\ref{chap:The Regulatory and Professional Framework}, we turn from this assessment to the next domain in the Prudentia framework.}
